Mental-health support dialogue is often evaluated as if the central problem were producing a warm and fluent reply. That target is too weak. Consider a student who describes exam catastrophizing, shame, or isolation. A useful support system should not merely reassure them; it should know whether the next turn should validate distress, elicit the situation--thought--emotion link, name a cognitive pattern, ask a Socratic question, plan a small action, guide grounding, or move to a safety boundary. These choices are process decisions, and they are easy for a fluent language model to obscure.

This issue is sharper in Vietnamese student support. Help-seeking conversations may combine academic pressure, family expectations, stigma, reluctance to seek formal care, and loneliness. Existing empathy-oriented benchmarks do not test whether a system can hold a therapy route across turns, pace CBT/BA/MBI interventions, decide when peer-like support should be silent, or expose enough state for post-hoc audit. Two responses may sound equally supportive while failing for different reasons: premature action advice, wrong modality, repeated peer commentary, missed crisis handling, or generic reassurance.

We study this problem as \emph{auditable process control}. \trace\ maintains a dynamic blackboard formulation that records the evolving dialogue state: presenting concern, distress, automatic thoughts, route evidence, stage milestones, delivered therapist interventions, peer boundary signals, safety flags, and benchmark metadata. A clinical assessor updates this state; deterministic controllers constrain stage progression; optional peer agents contribute only through a gated Yalom-factor policy; the therapist orchestrator decides how to use or discard those drafts; validator and safety modules repair or replace responses when needed.

The paper is a system and evaluation contribution, not a clinical efficacy claim. Its core argument is that mental-health dialogue systems should be evaluated not only by the surface quality of the final message, but also by the process decisions that led to it. We make four contributions:
\begin{itemize}[leftmargin=*]
    \item a dynamic blackboard formulation for Vietnamese student support dialogue, tracking distress, thoughts, route evidence, stage milestones, safety state, and intervention history;
    \item \trace, a process-controlled response architecture with assessor, therapist orchestrator, optional peer agents, validator, and safety critic;
    \item a Yalom-factor peer-gating mechanism that decides when peer support should appear, stay silent, or be replaced by therapist-only response;
    \item VSM, a benchmark protocol for route fidelity, stage fidelity, technique fidelity, safety, peer dynamics, latency, and ablation analysis.
\end{itemize}
